Este projeto investiga a atualização dinâmica de \textit{component trees}, 
especialmente a \textit{max tree}, para lidar com alterações
locais em imagens sem reconstruir toda a estrutura; 
a morfologia matemática fornece o alicerce teórico para essa
questão, baseada em conectividade, limiares e componentes
conexas. O problema central é que mudanças no nível de cinza 
podem fragmentar, fundir ou deslocar nós da árvore, tornando 
a representação anterior inconsistente e exigindo reprocessamento 
completo. Assim, o estudo propõe formalizar e desenvolver 
a operação de \textit{split}, que divide um nó em 
múltiplos filhos quando uma componente se fragmenta, 
buscando eficiência, correção e atualizações locais. 
 
Os objetivos incluem responder como a árvore deve ser 
modificada, quais estruturas de dados permitem manter o 
custo baixo, como recalcular atributos incrementalmente e 
qual ganho prático se obtém comparada a reconstrução completa.